\documentclass[11pt]{article}
\usepackage[a4paper,margin=23mm]{geometry}
\usepackage[T1]{fontenc}
\usepackage{lmodern,amsmath,amssymb,amsthm,mathtools,microtype,booktabs}
\usepackage[colorlinks=true,linkcolor=blue,citecolor=blue,urlcolor=blue]{hyperref}
\newcommand{\EE}{\mathbb E}
\newcommand{\tr}{\operatorname{tr}}
\newcommand{\WW}{\mathcal W}
\newcommand{\RR}{\mathcal R}
\newcommand{\LL}{\mathcal L}
\newcommand{\NN}{\mathcal N}
\newcommand{\OO}{\mathsf O_q}
\newcommand{\HH}{H^{(q,C)}}
\newcommand{\crs}{\operatorname{cr}}
\newcommand{\inv}{\operatorname{inv}}
\newcommand{\Reg}{\operatorname{Reg}}
\newtheorem{proposition}{Proposition}
\newtheorem{lemma}{Lemma}
\title{A constructive neutral memory kernel\\
\large Ordered Wick descendants, energy regression, and exact static removal}
\author{Alan Day \and Chengheng\thanks{AI research assistant (OpenAI).}}
\date{7 November 2026\quad Local research draft}
\begin{document}\maketitle
\begin{abstract}
We extend the preceding rank-marked chord rule to all ordered
Gaussian-Wick descendants in the neutral memory field. Internal field
pairings give an exact rule for Gaussian-chaos projections. A cut-diagram
argument identifies deterministic energy regression and cancels all
diagrams without a bridge between the two fields. Every surviving diagram
has exponential suppression in the external Hermite degrees. Consequently
each fixed descendant, after regression, has an absolutely continuous
two-energy spectrum. This permits the limit of exact samplewise energy
pinching. The previously proved uniform time-series tails then give a
convergent construction of the full neutral finite-time covariance.
The charged defect and its cross covariance remain to be computed.
\end{abstract}

\section{Setting and the remaining neutral problem}
Use standard CAR, overlapping complex Gaussian $p$-body couplings,
the same original couplings in all charge sectors, and the normalized
annealed full trace at $\beta=\mu=0$. Take even $N$, $p\to\infty$,
$\lambda_N=p^2/N\to\lambda>0$, fixed $J_s$, $\hbar=1$, and string time.
Set $C=J_s^2/\lambda$, $q=e^{-4\lambda}\in(0,1)$ and
\begin{equation}
 Z_j=\WW_j(H^j),\quad U_j=(jp-\RR)Z_j,
 \quad P_j(q)=[j]_q!,\quad \mathcal D_j=j^2-4q\partial_q.
\end{equation}
Here $\WW_j$ is ordinary coupling-Gaussian chaos projection.
All polynomial degrees are fixed before the double-scale limit.
The previous stage \cite{Size} replaces $U_j$ strongly by
$(jp-\LL)Z_j$, where $\LL$ is half the Clifford length, and gives
their two-mark chord moments. The preceding series \cite{Series}
reduces the neutral memory to these fields and forced lower-chaos
descendants, with uniform summable tails. We now compute those
descendants and their static part. No independent-probe substitution
is made.

\newpage
\section{Primitive descendants and the Gaussian-chaos label}
Write
\begin{equation}
 T_{v,w}=U_v\diamond Z_w,\qquad
 F_{r,v,w}=H^r T_{v,w},\quad v\ge1,\ w,r\ge0.
\end{equation}
The diamond is the ordinary Gaussian Wick product with ordered
matrix coefficients. Its entire $v+w$ source block forbids internal
Gaussian pairings; the physical size mark acts only on its first
$v$ slots. The $r$ ordinary slots may contract among themselves or
with the source. Taking an adjoint reverses all slots, so a prefix
mark becomes a suffix mark. These two types of block must be stored
separately.

For two primitive fields separated by $H^a,H^b$, sum over complete
pairings $\pi$ obeying the source-block restrictions. If $L_F=r+v+w$
and $i_F(\pi)$ pairs have both ends in the whole field interval, put
\begin{equation}
 d_F(\pi)=L_F-2i_F(\pi).\label{eq:degree}
\end{equation}
This is the degree of the field's uncontracted Gaussian Wick tensor.
An insertion of $\WW_dF$ keeps only $d_F=d$; an insertion of
$\OO F=q^{\NN}F$ multiplies the diagram by $q^{d_F}$.
External energy slots are never counted in $L_F$.

Let $I_1,I_2$ be the chord sets incident to the two marked subblocks,
and use the preceding geometrical weight
\begin{equation}
 K(I)=|I|^2-4\crs(\pi|_I),\quad
 K_1=K(I_1),\ K_2=K(I_2),\ K_{12}=K(I_1\cap I_2).
\end{equation}
The zero-, one-, and two-mark coefficients are respectively
\begin{equation}
 1,\qquad \lambda K_1,\qquad
 \lambda K_{12}+\lambda^2K_1K_2,\label{eq:marks}
\end{equation}
all multiplied by $C^{L/2}q^{\crs(\pi)}$, where $L$ is the total
number of raw Hamiltonian slots. Explicit scalar factors of $C$
in a descendant recurrence are retained in addition to this weight.
Odd $L$ gives zero. This is an algorithm for every fixed degree,
including chaos projections; it is not an assumption of spectral
$q$-Wick ordering at finite $N$.

The extension from rank to local size remains strong: right Wick
creation satisfies
\begin{equation}
 \|X_v\diamond Z_w\|_2\le
 \sqrt{(v+w)!/v!}\,M^{w/2}\|X_v\|_2,
 \qquad \sup_m C_m\le M.
\end{equation}
Fixed ordinary powers of $H$ are bounded on fixed Gaussian degree,
and chaos projections and $q^{\NN}$ are contractions. Thus only the
already proved two-field $L^2$ bridge is needed.

\newpage
\section{Energy regression from a cut field diagram}
Let $\HH_d$ denote the monic $q$-Hermite polynomial,
\begin{equation}
 \HH_{d+1}(x)=x\HH_d(x)-C[d]_q\HH_{d-1}(x),
 \qquad \|\HH_d\|_{L^2(\rho_q)}^2=C^dP_d(q),
\end{equation}
where $\rho_q$ is the energy law of variance $C$.
Define $\Reg F$ by its limiting inner products with energy
polynomials. For a pure source one obtains
\begin{equation}
 \boxed{\quad \Reg T_{v,w}=\kappa_v\HH_{v+w},\qquad
 \kappa_v=\lambda\frac{\mathcal D_vP_v(q)}{P_v(q)}.\quad}
 \label{eq:source-reg}
\end{equation}
Indeed, pair the source with a Wick block of length $v+w$.
All chords are bridges. Their permutations have weight
$q^{\inv\sigma}$. The Mahonian decomposition into the first $v$
positions, the remaining $w$ positions, and their shuffle shows
that the first group's internal permutation has its own
$q^{\inv}$ law. Its conditional average of $K_1$ is
$\mathcal D_vP_v/P_v$, independent of the selected endpoint set.
The suffix case follows by the same decomposition.

There is a finite formula for every projected primitive. Fix a
partial matching $\gamma$ internal to $F_{r,v,w}$, respecting
the source restriction. Its remaining $d=L_F-2|\gamma|$ slots are
Gaussian Wick ordered. Cut open the field interval and number all
their external partners in circular order. Crossings of internal
chords with one another and with open legs give a number $c_\gamma$
independent of the external assignment. Crossings between open
legs have count $\inv\sigma$ for a suitable order convention.
The mark has a value $K_\gamma(\sigma)$ depending only on this
finite configuration. Thus its regression coefficient is
\begin{equation}
 g_{\gamma,d}=C^{|\gamma|}q^{c_\gamma}\lambda
 \frac{\sum_{\sigma\in S_d}q^{\inv\sigma}K_\gamma(\sigma)}{P_d(q)}.
 \label{eq:cut-reg}
\end{equation}
Summing over $\gamma$ and the field's scalar coefficients gives
$\Reg(\WW_dF)=g_{F,d}\HH_d$. The case $d=0$ is included.
Formula \eqref{eq:degree} shows that $q^{\NN}$ simply multiplies
this group by $q^d$.

This derivation can also be checked by pairing with one external
$\HH_d$: the same permutation sum is exactly
$\langle F,\HH_d(H)\rangle/(C^dP_d)$.
The leading $Z_d-\HH_d(H)\to0$ bridge is used only in such linear
inner products. It is not used to infer continuity of arbitrary
products of marked fields.

\newpage
\section{No-bridge cancellation and the external Hermite bound}
Consider two projected primitives after subtracting their energy
regressions. Fix their internal matchings $\gamma,\gamma'$.
First sum diagrams having no chord directly joining the two field
intervals. Then $I_1\cap I_2=\varnothing$, so $K_{12}=0$.
With all external partner positions fixed, the two open-endpoint
permutations are independent. Reordering endpoints within one
field does not change crossings with any chord having both endpoints
outside that field. Therefore the two one-mark permutation sums
factor exactly as \eqref{eq:cut-reg}.

Subtracting regression cancels this entire no-bridge group,
including all internal partial matchings. The statement remains
true with $q^{\NN}$ or a fixed chaos projection, whose coefficients
are constant on each internal-matching group. An individual
no-bridge diagram need not vanish.

\begin{proposition}\label{prop:HS}
Let $R_F=F-\Reg F(H)$ and $R_G=G-\Reg G(H)$, where $F,G$ are
finite combinations of the primitives, their chaos projections,
and $q^{\NN}$ images. Their limiting two-energy measure has an
$L^2(\rho_q\otimes\rho_q)$ density. In particular it has no mass
on the diagonal $E=E'$.
\end{proposition}
\noindent\emph{Proof.}
Use the normalized external basis
$p_n=\HH_n/\sqrt{C^nP_n}$ and set
\begin{equation}
 M_{mn}=\lim\EE\tr[R_F^\dagger p_m(H)R_Gp_n(H)].
\end{equation}
The external $q$-Hermite polynomials implement the standard
multiblock Wick cancellations \cite{EP}: no pair remains internal
to either energy block. A deleted energy-internal chord belongs to
neither marked incident set; deleting it changes no crossing between
the remaining chords. Thus the marked weight is preserved as a
coefficient in this cancellation. This is a diagram identity.

Let $D$ bound the sum of the two raw field lengths. After the
no-bridge cancellation, there is at least one field-to-field bridge.
If $r$ chords connect the two energy blocks, endpoint counting gives
\begin{equation}
 |m-n|\le D,\qquad r\ge\min(m,n)-D.
\end{equation}
Every field bridge crosses every energy bridge, hence contributes
at least a factor $q^r$. Choosing the finite field attachments
costs at most a constant times $(m+n+1)^D$. The remaining energy
bridge permutations have total weight $P_r(q)$. Since
\begin{equation}
 \frac{P_r(q)}{\sqrt{P_m(q)P_n(q)}}\le(q;q)_\infty^{-1},
 \qquad r\le\min(m,n),
\end{equation}
the finite mark and regression coefficients give
\begin{equation}
 \boxed{|M_{mn}|\le K_D(m+n+1)^D
 q^{\max(0,\min(m,n)-D)}\,1_{|m-n|\le D}.}\label{eq:HS}
\end{equation}
Here $K_D$ depends on the fixed fields, $C,\lambda,q$, never on
$m,n$. In particular $\sum_{m,n}|M_{mn}|^2<\infty$.\hfill$\square$

\newpage
\section{From the bound to the spectrum and exact pinching}
For a finite field list define the positive matrix measures by
\begin{equation}
 \nu_{jk,N}(A,B)=\EE\tr[R_{j,N}^\dagger P_H(B)R_{k,N}P_H(A)].
 \label{eq:measure}
\end{equation}
Their fixed monomial moments follow from the finite diagram rule.
Absolute marked weights are bounded by fixed constants times
ordinary $q$-Gaussian moments. Positivity, higher-moment uniform
integrability and spectral Cauchy--Schwarz give weak convergence
to a unique measure supported on
\begin{equation}
 [-R_H,R_H]^2,\qquad R_H=2\sqrt{C/(1-q)}.
\end{equation}
As in the preceding stage, this concerns fixed field degrees and
all fixed external degrees; no uniform finite-$N$ estimate in the
external degree is asserted.

The bound \eqref{eq:HS} gives the density constructively:
\begin{equation}
 f_{jk}(E,E')=\sum_{m,n\ge0}M^{jk}_{mn}p_n(E)p_m(E'),
 \qquad d\nu_{jk}=f_{jk}\,d\rho_q(E)d\rho_q(E').\label{eq:density}
\end{equation}
The series converges in $L^2$. Its polynomial moments agree with
the original measure, and polynomials are dense on the compact
support, proving equality of measures. The matrix density is
positive almost everywhere. This uses the Hilbert--Schmidt
criterion also underlying \cite{BM}, not an unproved identification
of these common-coupling fields with independent $q$-Fock probes.
The diffuse measure $\rho_q\otimes\rho_q$ assigns zero mass to
$\Delta=\{E=E'\}$.

For a residual field with this limit, exact pinching satisfies
\begin{equation}
 \|\mathcal P_HR_N\|_2^2=\nu_{RR,N}(\Delta),\qquad
 \mathcal P_HX=\sum_eP_eXP_e.
\end{equation}
The diagonal is closed. Weak convergence of the finite positive
measures and Portmanteau therefore imply
\begin{equation}
 \limsup_N\|\mathcal P_HR_N\|_2^2
 \le\nu_{RR}(\Delta)=0.\label{eq:pinching}
\end{equation}
This retains entire degenerate energy blocks and proves annealed
$L^2$ convergence; it is not an interchange of an infinite-time
average with the double-scale limit.

An explicit coefficient-tail bound follows from \eqref{eq:HS}:
\begin{equation}
 \sum_{\min(m,n)\ge L}|M_{mn}|^2
 \le2(2D+1)K_D^2\sum_{n\ge L}(2n+D+1)^{2D}q^{2(n-D)}
 \quad(L\ge D).
\end{equation}
The scalar series has a geometric upper bound once its successive
ratio is below one. Also $(q;q)_\infty^{-1}\le
\exp[q/(1-q)^2]$ supplies a finite explicit normalization bound.

\newpage
\section{The forced descendants and their full static polynomials}
Suppress finite-$N$ subscripts after taking each fixed-degree limit.
The source and descendant recurrence are
\begin{align}
 r_j&=4\lambda C\sum_{v<j}q^vT_{v,j-1-v},\\
 \eta_{j+1}&=H\eta_j-C[j]_q\eta_{j-1}-r_j,
 \qquad \eta_0=\eta_1=\eta_2=0,
\end{align}
with $T_{0,w}=0$. Each $\eta_j$ is a finite linear combination of
$H^rT_{v,w}$, so all its diagrams are covered above.
If $g_j=\Reg\eta_j$, energy regression commutes with multiplying
by $H$, giving the scalar recurrence
\begin{equation}
 \boxed{g_{j+1}(x)=xg_j(x)-C[j]_qg_{j-1}(x)
 -4\lambda C\left(\sum_{v<j}q^v\kappa_v\right)\HH_{j-1}(x).}
 \label{eq:g}
\end{equation}
Again $g_0=g_1=g_2=0$. Write
$g_j=\sum_{d\le j-2}g_{j,d}\HH_d$. Self-adjointness of Gaussian
chaos projection, or the cut rule, gives
\begin{equation}
 \Reg(\OO\eta_j)=\sum_dq^dg_{j,d}\HH_d.
\end{equation}
The unpinched coefficients of the neutral series are
\begin{align}
 V_j&=4\lambda q^jU_j+(q^j-\OO)\eta_j,\\
 v_j(x):=\Reg V_j
 &=4\lambda q^j\kappa_j\HH_j(x)
    +\sum_{d\le j-2}(q^j-q^d)g_{j,d}\HH_d(x),\label{eq:v}\\
 R_j&=V_j-v_j(H).
\end{align}
Every $R_j$ has the absolutely continuous spectrum
\eqref{eq:density}, and its exact pinching vanishes in the limit.

For calibration,
\begin{align}
 g_3&=-4\lambda^2Cq\HH_1,\\
 g_4&=-\frac{8\lambda^2Cq(1+3q)}{1+q}\HH_2
       -4\lambda^2C^2q.
\end{align}
The scalar in $g_4$ agrees with the previously retained chaos-zero
term. No conserved term is removed before the recurrence or
$q^{\NN}$ operation. Moreover
\begin{equation}
 R_1=4\lambda q(U_1-\lambda H),\qquad
 d\nu_{11}=16\lambda^3Cq^2\,d\mu_q,
\end{equation}
where $\mu_q$ is the normalized Mehler measure of weight $q$.
The neutral memory slope $2CR_1$ therefore reproduces the previous
force spectrum $64\lambda^3C^3q^2\mu_q$.

\newpage
\section{The full neutral finite-time covariance}
Let
\begin{equation}
 c_j(s)=\frac{\int e^{isE}\HH_j(E)\,d\rho_q(E)}{C^jP_j(q)}.
\end{equation}
These are the limiting Jacobi coefficients in \cite{Series}.
That stage proved, with finite parameters retained until the limit,
\begin{equation}
 B_{m,N}(s)=\frac{m\epsilon_mC_m}{p}(1-\mathcal P_h)
 e^{-ihs/2}\sum_{j\ge0}c_{j,m}(\alpha_ms)V_{j,m,N}
 e^{-ihs/2}+o_{L^2}(1),
\end{equation}
uniformly on compact real memory times in the full charge sum.
Here $\alpha_m=(1-p/m)_+$, $\epsilon_m=1-\alpha_m^2$;
on central windows $\alpha_m\to1$, $m\epsilon_m/p\to2$ and
$C_m\to C$. The auxiliary Jacobi entries converge and are uniformly
bounded, so their operators converge strongly. Their exponentials
converge strongly uniformly on compact $s$-intervals; this gives
each $c_{j,m}(s)\to c_j(s)$. Continuity handles $\alpha_m\to1$.
The fixed-degree field coefficients converge to \eqref{eq:v}.
The tails are bounded uniformly by
an explicit summable envelope of the form
\begin{equation}
 \|c_{j,m}(s)V_{j,m,N}\|_2
 \le K_S(j+1)^2X_S^j/\sqrt{j!}.
\end{equation}

In the limiting joint Gram space, $\Reg$ is the orthogonal
projection onto the energy-polynomial closure. Hence
$\|R_j\|_2^2=\lim_N\|V_{j,N}\|_2^2-\|v_j\|_{L^2(\rho_q)}^2$
is at most $\lim_N\|V_{j,N}\|_2^2$, so the same coefficient
envelope controls the residual series.
For a fixed cutoff $J$, subtract $v_j(H)$ and apply
\eqref{eq:pinching} to every coefficient. The removed polynomials
commute exactly with $H$. Then remove $J$ using the above tail and
the contraction property of pinching. This proves the neutral
covariance limit
\begin{align}
 &\lim_N\EE\tr[B_N(s)^\dagger e^{iTH}B_N(u)e^{-iTH}]
 \notag\\
 &=4C^2\sum_{j,k\ge0}\overline{c_j(s)}c_k(u)
 \int e^{\frac i2(s-u)(E+E')+iT(E'-E)}\,d\nu_{jk}(E,E').
 \label{eq:kernel}
\end{align}
The sum is absolutely convergent, uniformly for $s,u$ in compact
sets and every real $T$: spectral Cauchy--Schwarz bounds each
measure's total variation by $\|R_j\|_2\|R_k\|_2$.
This all-$T$ statement concerns the series tail. The spectral
$N\to\infty$ limit is for each fixed $T$, not uniformly on the
whole real axis.
All $M^{jk}_{mn}$ in \eqref{eq:density} are supplied by the finite
diagram algorithm, so \eqref{eq:kernel} is a convergent constructive
answer, not a finite list of Taylor coefficients.

The spectral measures of the truncated time fields converge in
total variation as $J\to\infty$. Absolute continuity is closed
under this convergence, so the full neutral residual has no static
spectral mass. An $L^2$ density for the infinite series is not
required or asserted. Central uniformity and the previously proved
binomial charge tails give the full-trace result. No finite-temperature,
growing-time, or simultaneous $\lambda\to0$ limit is included.

\newpage
\section{Verification and the next physical calculation}
The checks accompanying this stage retain complete finite CAR
operators for Gaussian-chaos bookkeeping and separately test the
limiting chord weights. They compare finite Hermite expansions
with the internal-pair rule \eqref{eq:degree}; these are algebraic
checks, not finite-$N$ rank--size equalities. Independent permutation
enumeration tests \eqref{eq:source-reg}, and explicit no-bridge
groups test the cancellation used in Proposition~\ref{prop:HS}.
The descendant driver passed 42 complete finite comparisons
(maximum error $6.67\cdot10^{-15}$), 12 earlier-moment calibrations,
10 source regressions and 6 nonempty no-bridge cancellations.
It saves 59 signed moment polynomials and 9 chaos norms;
31 negative coefficients are retained. Finite normalized CAR
comparisons use floating arithmetic; formal polynomials are exact.

An independent driver enumerates 5,913 permutations through degree
seven, checks 56 prefix/suffix integer-polynomial identities, and
retains 12,316 negative mark values. Seven eight-slot no-bridge
layouts give 168 allowed pairings and 38 partner groups, all
with exact factorization. Their 120 nonzero centered single-diagram
weights include 60 negative weights. This last check uses formal
rational values $q=1/3$, $\lambda=1/4$ for polynomial identities,
not a physical parameter pair satisfying $q=e^{-4\lambda}$.
Both scientific drivers passed their first runs. Parameters,
versions and input hashes are recorded with the outputs.

\paragraph{What has changed.}
The neutral branch now has fixed-degree joint moments for all
forced descendants, their full deterministic static polynomials,
an all-external-degree spectral bound, and a controlled finite-time
covariance after exact energy pinching. The proof of the spectral
bound uses the arbitrary-degree cut argument and bridge counting;
finite enumerations alone would not establish absolute continuity.

\paragraph{Literature and attribution.}
Ordinary Gaussian Wick expansion, the multiblock $q$-Wick formulas
of Effros--Popa \cite{EP}, and the Hilbert--Schmidt strategy for
generator measures in Bikram--Mukherjee \cite{BM} are established
tools. Their use here requires the physical rank-to-size bridge,
the common-coupling local marks, and the cut-field cancellation
above. We do not infer a $q$-Gaussian ambient algebra for the
fluctuation fields solely from the energy law being $q$-Gaussian.

\paragraph{Continuation.}
The actual memory field remains $\mathsf F=A+B+o_{L^2}(1)$,
where $A$ is the charged transfer defect. The next calculation
must put its forward and reverse sources into the same common-$J$
diagram system and determine $AA$, $AB$, and $BA$. The nonzero
low-order mixed measure from \cite{Two} is a required calibration;
the neutral answer does not permit deleting these cross terms.
Once the full joint kernel is assembled, we can return to the
Susskind bulk mass/position dictionary and interaction observables.
The complete de Sitter dual and finite-$G$ four-point dynamics are
not established by this microscopic neutral calculation.

All artifacts remain local in this stage. There is no Git commit,
push, archive, or external issue reply.

\newpage
\begin{thebibliography}{9}\small
\bibitem{EP} E. G. Effros and M. Popa, \emph{Feynman Diagrams and Wick
Products Associated with $q$-Fock Space},
\href{https://arxiv.org/abs/math/0303045}{arXiv:math/0303045},
Theorems 3.2--3.3.
\bibitem{BM} P. Bikram and K. Mukherjee, \emph{Generator masas in
$q$-deformed Araki-Woods von Neumann algebras and factoriality},
\href{https://arxiv.org/abs/1606.04752}{arXiv:1606.04752v2},
Section 5, especially Theorems 5.3--5.4.
\bibitem{Size} A. Day and Chengheng, \emph{Rank marks as local
fermion-size defects}, local stage, 7 Nov. 2026,
\texttt{susskind-memory-size-bridge}.
\bibitem{Series} A. Day and Chengheng, \emph{A rank-marked series for
the neutral memory field}, local stage, 7 Nov. 2026,
\texttt{susskind-memory-neutral-marked-series}.
\bibitem{Two} A. Day and Chengheng, \emph{Four-field regression and a
two-component memory field}, local stage, 7 Nov. 2026,
\texttt{susskind-memory-two-component}.
\end{thebibliography}
\end{document}
